\documentclass[10pt,a4paper]{amsart}
\usepackage[margin=19mm]{geometry}
\usepackage{amsmath,amssymb,mathtools}
\usepackage[scr=rsfs,cal=euler]{mathalfa}
\usepackage{xfrac}
\usepackage[unicode,hidelinks,bookmarks=false]{hyperref}
\newcommand{\ie}{i.e.\ }
\newcommand{\cf}{cf.\ }
\newcommand{\resp}{respectively\ }
\newcommand{\Subsection}{\subsection}
\providecommand{\nobreakdash}{\nobreak\hbox{-}}
\setlength{\emergencystretch}{2em}
\multlinegap0pt
\usepackage{amssymb}
\usepackage{mathtools}
\usepackage{tcolorbox}
\usepackage{algorithm}\usepackage{algpseudocode}
\usepackage{xcolor}
\usepackage[shortlabels]{enumitem}
\usepackage{tikz}
\newtheorem{theorem}{Theorem}
\title{On the convexity of Sum of Separable and Biquadratic Forms}
\author{Shaon Naskar, Sujeet Kumar Singh}
\date{Source: arXiv:2607.23476v1 (CC BY 4.0)}
\begin{document}
\maketitle

\noindent\emph{Source and licence.} On the convexity of Sum of Separable and Biquadratic Forms. Version arXiv:2607.23476v1 (CC BY 4.0), distributed by arXiv under the Creative Commons Attribution 4.0 International licence, http://creativecommons.org/licenses/by/4.0/, as recorded in the arXiv licence metadata for this exact version and shown on the abstract page https://arxiv.org/abs/2607.23476v1. Full licence notice: \texttt{licenses/arxiv-2607.23476-CC-BY-4.0.txt}.\par\smallskip

\noindent\textit{Editorial scope.} Counted unit: the article's theorem theorem (source0.tex lines 347--349). The article's complete original proof of it is delivered with it (source0.tex lines 351--429, one \texttt{proof} environment of 79 lines). No mathematics is written, repaired or re-derived: the statement and the proof are the article's own lines, in the article's own notation, and no retained line is substituted or re-flowed.  Retained uncounted as the source-local context the proof needs: lines 766--766.  Nothing was omitted from any retained span, so the package carries no omission record.  No retained line contains a citation, so the wrapper carries no bibliography and no bibliography entry.  No retained line contains a figure environment, an image-inclusion command or drawing code, so no image is delivered and the package declares no asset dependency.  Editorial formatting only, in addition: an AMS \texttt{amsart} preamble that restates the article's own declarations and macros used by the retained lines, namely the environments \texttt{theorem} on separate counters, together with the standard packages those lines need.  The retained environments print in this document as Theorem 1 in order of appearance, whereas the article prints the same items with the numbering of its own full document.  The complete native source of the article as distributed in the arXiv e-print for this exact version is bundled unchanged as \texttt{sources/206031/source0.tex}.
\section*{Appendix B}\label{Appendix:B}

\begin{theorem}
    The polynomial $f(x,y) = f_{sep}(x,y) + 0.1 B^*(x,y)$ is not sos-convex.
\end{theorem}

\begin{proof}

Suppose 
\[
p = z^T \nabla^2 f z = \sum_i p_i^2,
\]
where each $p_i(x,y)$ is a bilinear form. We prove that the polynomial $z^T \nabla^2 f z$ in 12 variables $(x_1,x_2,x_3,y_1,y_2,y_3,z_{x_1},z_{x_2},z_{x_3},z_{y_1},z_{y_2},z_{y_3})$ of degree $4$ is not sos. The explicit expression of $z^T \nabla^2 f z$ is given in Appendix B \ref{Appendix:B}.

Observe that the terms $x_i^2 z_{x_j}^{2}$, where $i\neq j \in \{1,2,3\}$, and $y_m^2 z_{y_n}^{2}$, where $m\neq n \in \{1,2,3\}$ are absent in $p$, which means they are also absent in $p_i^2$, and hence the terms are absent in $p_i$. Rewrite $p_i = s_i + t_i$, where $s_i$ contains only the terms $x_i z_{x_i}, y_i z_{y_i}$ for $i=1,2,3$, and $t_i$ contains only the terms $x_i z_{y_j}, y_iz_{x_j}$ for $i,j =1,2,3$.

Now, equating the coefficients of $p = \sum_i p_i^2 = \sum_i (s_i + t_i)^2$ leads to the following:

\[
    2s_i t_i =0
\]

\begin{align*}
    s_i^2 =\; &+ 0.4\Big[x_1z_{x_1} + y_2z_{y_2}\Big]^2
+ 0.4\Big[x_2z_{x_2} + y_3z_{y_3}\Big]^2
+ 0.4\Big[x_3z_{x_3} + y_1z_{y_1}\Big]^2 \\
&+ 0.4\Big[x_2z_{x_2} + y_1z_{y_1}\Big]^2
+ 0.4\Big[x_3z_{x_3} + y_2z_{y_2}\Big]^2
+ 0.4\Big[x_1z_{x_1} + y_3z_{y_3}\Big]^2 \\
&+ 0.4\Big[x_1z_{x_1} + y_2z_{y_2}\Big]^2
+ 0.4\Big[x_2z_{x_2} + y_3z_{y_3}\Big]^2
+ 0.4\Big[x_3z_{x_3} + y_1z_{y_1}\Big]^2 \\
&+ 0.4\Big[x_2z_{x_2} + y_1z_{y_1}\Big]^2
+ 0.4\Big[x_3z_{x_3} + y_2z_{y_2}\Big]^2
+ 0.4\Big[x_1z_{x_1} + y_3z_{y_3}\Big]^2 \\
&+ 0.2\Big[x_1z_{x_1} + y_1z_{y_1}\Big]^2
+ 0.2\Big[x_1z_{x_1} + y_2z_{y_2}\Big]^2
+ 0.2\Big[x_2z_{x_2} + y_2z_{y_2}\Big]^2 \\
&+ 0.2\Big[x_2z_{x_2} + y_3z_{y_3}\Big]^2
+ 0.2\Big[x_3z_{x_3} + y_3z_{y_3}\Big]^2
+ 0.2\Big[x_3z_{x_3} + y_1z_{y_1}\Big]^2 \\
&+ 0.2\Big[x_1z_{x_1} + y_1z_{y_1}\Big]^2
+ 0.2\Big[x_2z_{x_2} + y_1z_{y_1}\Big]^2
+ 0.2\Big[x_2z_{x_2} + y_2z_{y_2}\Big]^2 \\
&+ 0.2\Big[x_3z_{x_3} + y_2z_{y_2}\Big]^2
+ 0.2\Big[x_3z_{x_3} + y_3z_{y_3}\Big]^2
+ 0.2\Big[x_1z_{x_1} + y_3z_{y_3}\Big]^2 \\
&+ 9.6\Big(x_1^2z_{x_1}^2 + x_2^2z_{x_2}^2 + x_3^2z_{x_3}^2
+ y_1^2z_{y_1}^2 + y_2^2z_{y_2}^2 + y_3^2z_{y_3}^2\Big).
    \end{align*}

\begin{align*}
     t_i^2 =\; & 0.2\Big[(y_1 - y_2)z_{x_1} + x_1(z_{y_1} - z_{y_2})\Big]^2 \\
&+ 0.2\Big[(y_2 - y_3)z_{x_2} + x_2(z_{y_2} - z_{y_3})\Big]^2 \\
&+ 0.2\Big[(y_3 - y_1)z_{x_3} + x_3(z_{y_3} - z_{y_1})\Big]^2 \\
&+ 0.2\Big[(x_1 - x_2)z_{y_1} + y_1(z_{x_1} - z_{x_2})\Big]^2 \\
&+ 0.2\Big[(x_2 - x_3)z_{y_2} + y_2(z_{x_2} - z_{x_3})\Big]^2 \\
&+ 0.2\Big[(x_3 - x_1)z_{y_3} + y_3(z_{x_3} - z_{x_1})\Big]^2 \\
&+ 0.4y_2^2z_{x_1}^2 + 0.4x_1^2z_{y_2}^2
+ 0.4y_3^2z_{x_2}^2 + 0.4x_2^2z_{y_3}^2 \\
&+ 0.4y_1^2z_{x_3}^2 + 0.4x_3^2z_{y_1}^2
+ 0.4y_1^2z_{x_2}^2 + 0.4x_2^2z_{y_1}^2 \\
&+ 0.4y_2^2z_{x_3}^2 + 0.4x_3^2z_{y_2}^2
+ 0.4y_3^2z_{x_1}^2 + 0.4x_1^2z_{y_3}^2 \\
&- 0.4x_1y_2z_{x_1}z_{y_1}
- 0.4x_1y_1z_{x_1}z_{y_2}
- 0.4x_2y_3z_{x_2}z_{y_2}
- 0.4x_2y_2z_{x_2}z_{y_3} \\
&- 0.4x_3y_1z_{x_3}z_{y_3}
- 0.4x_3y_3z_{x_3}z_{y_1}
- 0.4x_1y_1z_{x_2}z_{y_1}
- 0.4x_2y_1z_{x_1}z_{y_1} \\
&- 0.4x_2y_2z_{x_3}z_{y_2}
- 0.4x_3y_2z_{x_2}z_{y_2}
- 0.4x_3y_3z_{x_1}z_{y_3}
- 0.4x_1y_3z_{x_3}z_{y_3}.
\end{align*}
The last equality is absurd, as the right hand side takes a negative value when evaluated at:
\[
(x_1 =1, x_2=0, x_3=2, y_1=1, y_2=0, y_3=-2, z_{x_1}=-1, z_{x_2}=0, z_{x_3}=2, z_{y_1}=-1, z_{y_2}=0, \\z_{y_3}=2).
\]

Therefore, we have established that $p=z^T \nabla^2 f z$ cannot be expressed as the sum of squares of bilinear forms.

\end{proof}

\end{document}
